% \documentclass[aspectratio=169]{beamer}
\documentclass[aspectratio=169, handout]{beamer}
% Add handout option to ignore pauses

% From Jeff E
\usepackage{algo}
\usepackage{sigmastyle}

\setbeamertemplate{blocks}[rounded]
\setbeamercovered{transparent}
\usecolortheme{orchid}
\usetheme{sigma}

\graphicspath{{images/}}

\title{Generating a random sample}
\subtitle{Introduction to Discrete Probability}

\author{Ian Chen}
\date{}

% \institute{University of Illinois Urbana-Champaign}

% --------------------------------------------------------------------

% Begin document
\begin{document}

\begin{frame}
    \titlepage
\end{frame}

\section{Introduction}

\begin{frame}{Samples}
    \only<1>{
        \begin{itemize}
            \item Why do we want a random sample?
            \item What is a random sample?
            \item How do we generate a random sample?
        \end{itemize}
    }
    \only<2->{
        \begin{itemize}
            \item What is a random sample?
                  \begin{itemize}
                      \item<3-> Should be \emph{smaller} than the dataset
                      \item<4-> Each \emph{item} should have equal chances of being in the sample
                      \item<5-> Each \emph{pair of items} should have equal chances
                      \item<5-> $\ldots$
                  \end{itemize}
        \end{itemize}
    }
\end{frame}

\begin{frame}{Streams}
    \begin{itemize}
        \item A stream is a sequence of items
              \begin{itemize}
                  \item<2-> You don't know how many items you have
                  \item<3-> You can only read each item once, in order
              \end{itemize}
        \item<4-> We want to generate a sample from a stream
    \end{itemize}
\end{frame}


\section{Permutations}
\frame{\tableofcontents[currentsection]}

\begin{frame}{Lots}
    \begin{algo}
        \textul{\textsc{DrawLots}(array $L$)}: \+
        \\ $n \gets \abs{L}$
        \\ for $i$ from $1 \to n$
        \+ \\ remove a random lot $x$ from $L$ \-
        \+ \\ $R[i] \gets x$ \-
        \\ return $R[1 \ldots n]$
    \end{algo}

    \vspace{\baselineskip}

    \only<2->{
        \begin{itemize}
            \item Correctness?
                  \only<3>{
                      \begin{itemize}
                          \item There are $n$ equally likely choices for $R[1]$
                          \item There are $n-1$ equally likely choics for $R[2]$
                          \item $\ldots$ There are $n!$ equally likely results for $R$
                      \end{itemize}
                  }
        \end{itemize}
    }
\end{frame}

\begin{frame}{Fisher-Yates}
    \begin{algo}
        \textul{\textsc{InsertionShuffle}(stream $S$)}: \+
        \\ Initialize $A$ as empty array
        \\ while $S$ is not empty
        \+ \\ Append next item in $S$ to $A$ \-
        \+ \\ swap $A[\abs{A}] \leftrightarrow A[\textsc{Random}(\abs{A})]$
    \end{algo}

    \vspace{\baselineskip}

    \only<2->{
        \begin{itemize}
            \item Correctness?
                  \only<3>{
                      \begin{itemize}
                          \item There are $1$ equally likely choices for $\textsc{Random}(1)$
                          \item There are $2$ equally likely choices for $\textsc{Random}(2)$
                          \item $\ldots$ There are $n!$ equally likely results for the algorithm
                      \end{itemize}
                  }
        \end{itemize}
    }
\end{frame}

\section{Samples}
\frame{\tableofcontents[currentsection]}

\begin{frame}{Permutations to samples}
    \only<1>{
        \begin{itemize}
            \item Interpret the \emph{key} of an item as the position in the permutation
            \item The sample is created from the $k$ smallest keys
        \end{itemize}
    }

    \only<2->{
        \begin{itemize}
            \item Normalize keys between $[0, 1]$
        \end{itemize}

        \begin{algo}
            \textul{$k$-\textsc{Min-Sample}(stream $S$, int $k$)}: \+
            \\ Initialize $R \gets \emptyset$ for output
            \\ while $S$ is not done
            \+ \\ $x \gets $ next item in $S$ \-
            \+ \\ $x.key \gets $ random number between $[0, 1]$ \-
            \+ \\ Update $R$ to contain the $k$ smallest keys \-
            \\ return $R$
        \end{algo}
    }

    \only<3->{
        \begin{itemize}
            \item Correctness?
                  \begin{itemize}
                      \item<4-> Same as discretized version!
                  \end{itemize}
        \end{itemize}
    }
\end{frame}

\begin{frame}{Making it faster}
    \only<1-3>{
        \begin{itemize}
            \item Most of the time we do not update the reservoir
            \item<2-> Could we only generate random numbers when we need to update the reservoir?
            \item<3-> We need to know how many items to skip between random number calls
        \end{itemize}
    }

    \only<4->{
        \begin{itemize}
            \item What is the \emph{probability} of skipping $0$ items? $1$? $j$?
                  \only<5>{
            \item Geometric! $\textbf{Pr}(\text{skip } j \text{ items}) = p(1-p)^j$
                  }
        \end{itemize}
    }
\end{frame}

\begin{frame}[allowframebreaks]{Reservoir Sampling}
    \only<1>{
        \begin{itemize}
            \item Generate numbers according to the distribution of skipping $j$ items
        \end{itemize}

        \begin{algo}
            \textul{\textsc{Reservoir-Sample}(stream $S$, int $k$)}: \+
            \\ Initialize $R \gets \emptyset, p \gets 1$
            \\ while $S$ is not done
            \+ \\ $x \gets $ next item in $S$ \-
            \+ \\ $x.key \gets $ random key between $[0, M]$ \-
            \+ \\ Update $R$ to contain the $k$ smallest keys \-
            \+ \\ Update $p$ to be largest key in $R$ \-
            \+ \\ Create $j$ with distribution $p(1-p)^j$ and skip next $j$ items \-
            \\ return $R$
        \end{algo}
    }
    \only<2->{
        \begin{itemize}
            \item<2-> What is the \emph{expected} number of random numbers generated?
                  \begin{itemize}
                      \item<3-> Let $n$ be number of elements in stream
                      \item<3-> Let $I_i$ denote whether we generate a random number for element $i$
                            \only<4->{
                      \item<4-> $N$ is the total number of generated numbers, to find $\textbf{Ex}(N)$,
                            \begin{align*}
                                \textbf{Ex}(N)
                                 & \uncover<4->{= \textbf{Ex}( \sum_{i=1}^n I_i )}
                                \uncover<5->{= \sum_{i=1}^n \textbf{Ex}( I_i )}    \\
                                 & \uncover<6->{= k + \sum_{i=k+1}^n \frac{k}{i}}  \\
                                 & \uncover<7->{< k(1 + \ln(n/k)) + 1}
                            \end{align*}
                            \vspace{-2\baselineskip}
                      \item<8-> Much better than $n$
                            }
                  \end{itemize}
        \end{itemize}
    }
\end{frame}


\section{What's next?}
\frame{\tableofcontents[currentsection]}

% for some reason, allowframebreaks does't work well with this
\begin{frame}{Conclusions I}
    \begin{itemize}
        \item<1-> Generating geometric random variables
              \begin{itemize}
                  \item If $X \sim \text{Geom}(p)$, $U \sim \text{Uniform}(0, 1)$, then
                        $$
                            \textbf{Pr}(X < x) = \textbf{Pr}(\floor{\frac{\ln(U)}{\ln(1 - p)}} < x)
                        $$
              \end{itemize}
        \item<2-> Weighted sampling
              \begin{itemize}
                  \item If $X_i \overset{indep}{\sim} \text{Exp}(\lambda_i)$, then for $i \ne j$,
                        $$
                            \textbf{Pr}(X_i < X_j) = \frac{\lambda_i}{\lambda_i + \lambda_j}
                        $$
                        \vspace{-\baselineskip}
                  \item Generate keys according to exponential distribution
                  \item The probability of being in the sample is proportional to the weight
              \end{itemize}
    \end{itemize}
\end{frame}

\begin{frame}{Conclusions II}
    \begin{itemize}
        \item<1-> Getting rid of the heap
              \begin{itemize}
                  \item The position of the largest key is distributed uniformly in the array
                  \item Pick an element at random as the max!
                  \item Simulates the heap
              \end{itemize}
        \item<2-> Sketching, or making faster algorithms with random samples
              \begin{itemize}
                  \item Median selection
                  \item Cardinality estimation
                  \item Linear queries
              \end{itemize}
    \end{itemize}

    % \item<3> Online linear queries
    %       \begin{itemize}
    %           \item Suppose we wanted to calculate the $\sum_{i} h(i) A[i]$ for online $h$
    %           \item We can take a sample, (do some small tricks to generate $\hat{h}(i)$, and do
    %                 $$
    %                     \sum_{i \in S} \hat{h}(i) A[i]
    %                 $$
    % \item If $X_i \overset{indep}{\sim} \text{Exp}(\lambda_i)$,
    %       $$
    %           \textbf{min}_i (X_i) \sim \text{Exp}(\sum_i \lambda_i)
    %       $$
    %       \vspace{-\baselineskip}
    % \item Calculate a sample, take the minimum key,
    % \end{itemize}
\end{frame}


\end{document}
